\section{The Pipeline}
\label{sec:pipeline}

Figure~\ref{fig:pipeline} shows the stages of the tool. Every input format is converted to one common
representation, the \emph{article model}: the sentences in reading order, the paragraph and section of each
sentence, the figures with their captions, and for each sentence the set of figures that it cites. Everything
after parsing works on this representation only.

\begin{figure}[h]
\centering
\begin{tikzpicture}[
  box/.style={draw=navy!40, fill=soft, rounded corners=2pt, align=center, font=\small,
              minimum height=1.25cm, text width=2.35cm},
  dark/.style={box, fill=navy, text=white, draw=navy},
  arr/.style={-{Stealth[length=5pt]}, thick, accent}]
\node[box] (in) {\textbf{Input}\\ XML, \LaTeX{},\\ or PDF};
\node[box, right=0.45cm of in] (parse) {\textbf{Parse}\\ JATS / GROBID /\\ own PDF parser};
\node[box, right=0.45cm of parse] (art) {\textbf{Article model}\\ sentences, sections,\\ captions};
\node[box, right=0.45cm of art] (exp) {\textbf{Explicit}\\ citation\\ grammar};
\node[dark, right=0.45cm of exp] (imp) {\textbf{Implicit}\\ rule \bof{} or\\ learned models};
\node[box, below=0.6cm of imp] (out) {\textbf{Output}\\ JSON + HTML\\ report};
\draw[arr] (in)--(parse); \draw[arr] (parse)--(art); \draw[arr] (art)--(exp); \draw[arr] (exp)--(imp);
\draw[arr] (imp)--(out);
\end{tikzpicture}
\caption{Stages of the tool.}
\label{fig:pipeline}
\end{figure}

Papers can be read from structured XML, from \LaTeX{} source, or from PDF through GROBID or through our own
PyMuPDF-based layout parser; Appendix~\ref{app:parsing} describes these routes.

\subsection{Explicit references}
Explicit references are found by a regular-expression grammar. It recognises the keyword in its common forms
(Fig., Figure, Figs., Figures), lists (``Figs.~2 and 3''), ranges (``Figs.~2--4''), sub-figure letters (``3B'',
``4A,B''), and, only when the surrounding context supports it, a bare panel reference such as ``4(a)''. A
mention is kept only if it refers to a figure that was detected in the paper. At the sentence level it reaches
an F1 of 0.93 on 28 arXiv papers and 0.85 on 177 bioRxiv papers that were not used during its development.

\subsection{The command-line tool}
\begin{Verbatim}[fontsize=\small,frame=single,framesep=6pt,rulecolor=\color{navy!30}]
python -m figrefs <pdf or directory> -o out/
\end{Verbatim}
For each paper the tool writes \texttt{result.json} (for every figure: the cropped image, the caption, the
explicit sentences, the implicit sentences and any mentions that could not be matched to a figure) and a
self-contained \texttt{report.html}; for a directory it also writes \texttt{all\_results.json}. With the rule
\bof{} the tool needs no neural model and takes about four seconds on a CPU for a 21-figure paper. The learned
models run through a separate inference script (Section~\ref{sec:cs}); on a laptop they take about three to
four minutes per paper, and less on a GPU.
